\begin{Exercise}[label={ex:de_morgan_tabella}]
    Data l'espressione logica $\text{!(A || !B) \&\& C}$:
    \begin{enumerate}
        \item Utilizzare le Leggi di De Morgan per semplificare la negazione della parentesi.
        \item Costruire la tabella di verità completa per l'espressione originale.
    \end{enumerate}
\end{Exercise}

\begin{Answer}[ref={ex:de_morgan_tabella}]
    \begin{enumerate}
        \item Semplificazione (Leggi di De Morgan):
        \begin{align*}
            \text{!(A || !B) \&\& C} &\equiv (\text{!A \&\& !!B) \&\& C} \quad \text{[De Morgan]} \\
            &\equiv !A \&\& B \&\& C \quad \text{[Doppia Negazione]}
        \end{align*}
        \item Tabella di verità:
            \[
            \begin{array}{|c|c|c|c|c|c|}
                \hline
                \text{A} & \text{B} & \text{C} & \text{!B} & \text{A || !B} & \text{!(A || !B) \&\& C} \\ \hline
                0 & 0 & 0 & 1 & 1 & 0 \\ \hline
                0 & 0 & 1 & 1 & 1 & 0 \\ \hline
                0 & 1 & 0 & 0 & 0 & 0 \\ \hline
                0 & 1 & 1 & 0 & 0 & 1 \\ \hline
                1 & 0 & 0 & 1 & 1 & 0 \\ \hline
                1 & 0 & 1 & 1 & 1 & 0 \\ \hline
                1 & 1 & 0 & 0 & 1 & 0 \\ \hline
                1 & 1 & 1 & 0 & 1 & 0 \\ \hline
        
            \end{array}
            \]
    \end{enumerate}
\end{Answer}